\documentclass[12pt]{article}
\usepackage{amsmath}
\usepackage{geometry}
\geometry{letterpaper, margin=1in}
\usepackage{fancyhdr}
\usepackage{setspace}
\usepackage{hyperref}

\pagestyle{fancy}
\fancyhf{}
\rhead{33-120 Science and Science Fiction}
\lhead{Sci-Fi Problem 1 (25 points)}
\setlength{\headheight}{15pt}

\begin{document}

\vspace{1cm}

\noindent
\textbf{Concept:} acceleration = change in velocity over change in time

\[
a = \frac{\Delta v}{\Delta t} = \frac{v_{\text{final}} - v_{\text{initial}}}{\Delta t}
\]

\vspace{0.5cm}

\noindent You will receive 5 points for an on time submission. The final 20
points are assigned below. In order to get full credit for the problems below,
be sure to:
\begin{itemize}
    \item explicitly state all assumptions made that are not given in the
    problem.
    \vspace{-10pt}
    \item show all of the steps in each calculation.
    \vspace{-10pt}
    \item include proper physical units for all results.
\end{itemize}


\noindent
\textbf{Problem 1. (10 points)} Imagine human space flight from Earth to Mars in
the following two stages.

\begin{enumerate}
    \item[a.] The spaceship starts at rest on the surface of the Earth and
    accelerates at \textbf{three times} the acceleration due to gravity (\(a =
    3g\), where \(g = 9.8 \, \text{m/s}^2\)), typical of actual space launches.
    Calculate how long (\(\Delta t\)) would it take for the spaceship to reach
    \textbf{escape velocity} – the speed sufficient to enter orbit around the
    Earth. (Hint: look in lecture slides for a discussion of \textbf{escape velocity} from Earth)
    
    \item[b.] Now that the spaceship has reached escape velocity and is
    successfully in orbit around the Earth, the most efficient way to get from
    Earth orbit to Mars is by a 
    \href{https://en.wikipedia.org/wiki/Hohmann_transfer_orbit}{Hohmann Transfer
    Orbit}.  The increase in speed needed to do this is a change in velocity,
    \(\Delta v\), of approximately 2.9 km/s. The necessary increase in velocity
    is achieved by burning the engines of the spaceship for a certain amount of
    time. If the “burn time” is 90 seconds, what acceleration will the
    astronauts experience?
\end{enumerate}

\vspace{1cm}

\noindent


\vspace{0.5cm}

\noindent
\textbf{Problem 2. (10 points)} Consider the film clip from The Martian shown in
class which depicts Mark Watney's launch from the surface of Mars, starting from
zero initial speed and accelerating to a final speed of 850 m/s. The change in
speed takes place over a time interval of about 40 seconds.

\begin{enumerate}
    \item[a.] Convert the final speed of the spacecraft from meters per second
    to miles per hour using an online source (be sure to include a reference) or
    by converting miles to meters and hours to seconds. 
    \item[b.] Calculate the acceleration Watney experiences and compare it to
    the acceleration due to gravity on Earth.
    \item[c.] Is this scene plausible? Does it make sense that Watney would have
    passed out based on the acceleration he felt?
\end{enumerate}

\vspace{0.5cm}

\noindent

\newpage
\vspace{1cm}

\noindent {\Large Rubric}

\noindent
\textbf{Problem 1. (10 points)} Imagine human space flight from Earth to Mars in
the following two stages.

\begin{enumerate}
    \item[a.] 5 pts total; -1 for not writing out the equation for $\Delta t$;
    -1 for incorrect escape velocity; -1 for incorrect plugging in of numbers to
    get value for $\Delta t$. They should get at least 2 points for trying.

    $\Delta t = \frac{\Delta v }{a} = \frac{v_{\rm{esc}}}{a} = \frac{11,171\rm{m/s}}{3\times9.8\rm{m/s^2}} = 380 \rm{s}$
    
    \item[b.] 5 pts total; -1 for not writing out the equation for acceleration;
    -1 for incorrect plugging in of numbers. They should get at least 3 points
    for trying. 
    $a = \frac{\Delta v}{\Delta t} = \frac{2900\rm{m/s}}{90\rm{s}} = 32.2 \rm{m/s^2}$

\end{enumerate}

\vspace{1cm}

\noindent


\vspace{0.5cm}

\noindent
\textbf{Problem 2. (10 points)} Consider the film clip from The Martian shown in
class which depicts Mark Watney's launch from the surface of Mars, starting from
zero initial speed and accelerating to a final speed of 850 m/s. The change in
speed takes place over a time interval of about 40 seconds.

\begin{enumerate}
    \item[a.] Convert the final speed of the spacecraft from meters per second
    to miles per hour using an online source (be sure to include a reference) or
    by converting miles to meters and hours to seconds. 

    850 m/s = 1901.4 miles per hour (Wolfram Alpha)
    
    2 points; -1 for incorrect value; 1 pt for trying. 


    \item[b.] Calculate the acceleration Watney experiences and compare it to
    the acceleration due to gravity on Earth. 
    
    $a=\frac{\Delta v}{\Delta t} = \frac{850\,\rm{m/s}}{40\,\rm{s}} = 21.25\rm{m/s^2}$
    $a/g = 21.25/9.8 = 2.17$

    4 points; -1 for not writing equation; -1 for incorrect value; -1 for no comparison
    \item[c.] 4 points; -1 for incorrect logic
    
    Students should make an argument about whether they think it is realistic for him to have passed out. 
    As long as it is reasonable, they can get points. (most, if not all, should get points)
\end{enumerate}


\noindent

\end{document}
